\section{Critical review, corrections, and limits of the conclusions}
\label{sec:review}

The pinned manuscript already contains important corrections to earlier
stages of this project. Its $S_4$ reduction is a theorem, its constant-one
Euler enclosure has a stated domain, and several low-index Lerch motion
heuristics have already been replaced by proved phase information.
Those corrections should survive integration. The present review found
no reason to retract those proved results. The changes proposed below
update open items that the new arguments settle and prevent stronger
interpretations that the arguments do not support.

\subsection{The Gaussian maximum now controls every Euler index}

The caution in the caption of the Gaussian-axis figure in
\code{05-subcritical-stieltjes.tex}, and the corresponding paragraph in
\code{10-discovery.tex}, distinguish a first-remainder maximum from a
universal Euler constant. That distinction was justified: knowing the
largest value of $-2g_{a,b}$ alone does not bound the later scaled tails.
Lemma~\ref{newEuler:lem:scalar} supplies the missing comparison. The
replacement should state Theorem~\ref{newEuler:thm:domination} and
Corollary~\ref{newEuler:cor:sharp}, including the strict positive-order
inequality and the limiting equality at $a=0$.

The existing counterexample to a universal constant one remains valid.
The improved theorem changes the constant to the optimal $C_*>1$;
it does not extend the old constant-one theorem with its hypotheses
removed. Similarly, the stronger bounds already proved on smaller
parameter regions should remain available.

\begin{proposition}[Scaled Euler errors need not decrease]
\label{prop:scaled-error-counterexample}
For the positive integer pair $(a,b)=(2,1)$,
\[
 R_2(2,1)-R_1(2,1)>\frac1{60}>0.
\]
Nevertheless $R_N(2,1)\to0$. Thus, even at integer orders, the largest
scaled error need not occur at $N=1$.
\end{proposition}
\begin{proof}
The first relevant coefficients are
$c_3=1/6$, $c_5=1/12$, and $c_7=1/20$. Hence
$d_1=-1/6$, $d_2=-1/4$, $d_3=-3/10$, and
\[
 E_1=0,\qquad E_2=-\frac1{24},\qquad E_4=-\frac{11}{120}.
\]
The strict one-sided enclosure gives $g_{2,1}<E_4$. Therefore
\[
 R_2-R_1=4E_2-2E_1-2g_{2,1}
       =-\frac16-2g_{2,1}>\frac1{60}.
\]
The limit follows from Theorem~\ref{newEuler:thm:domination}.
This is a finite symbolic counterexample, independent of a numerical
plot or a floating-point evaluation of $g_{2,1}$.
\end{proof}

For a fixed positive pair, the positive sequence $R_N$ tends to zero,
so its supremum is attained at at least one finite index. Determining
that index and its dependence on $(a,b)$ is a separate problem from
the sharp constant uniform in both orders.

\subsection{Uniformity must be attached to an approximation and a domain}

The former compact-ratio moment theorem should retain its original
statement. Theorem~\ref{thm:all-ratio} supplies an additional expansion
whose relative remainder is uniform on $0\le m\le n$. The two phases
place different parts of the Jacobian in the exponential, and their
saddles need not agree exactly at finite exponents. Asymptotic
equivalence does not justify identifying those saddles term by term.

The all-ratio theorem gives any prescribed algebraic order in $n^{-1}$.
It does not replace exponentially accurate fixed-$m$ residue expansions,
or their least-term remainder estimates, by a bound of the same strength.
A future matched theorem must state the approximation and error scale
explicitly in each overlap.

For Lerch zeros, the new threshold $K_n$ proves completeness by sign
transfer and a multiplicity bound. A numerical scan would not suffice.
The threshold should always be described as sufficient. At index two
the source already proves saturation for every $k\ge1$, much stronger
than either of our general-purpose quantitative bounds. The new result
is useful because it applies to arbitrary $n$, arbitrary branches, and
the full closed deformation interval simultaneously.

\subsection{Formal coordinates and numerical periods}

The quotient in Corollary~\ref{oi:cor:basket} is defined by the
specified $m$ shuffle rows. Its dimension is exactly $m$ by an invertible
minor. Evaluating at $i$, or imposing additional functional equations,
may create further relations. Thus neither this dimension nor the Smith
form proves linear independence of the corresponding real constants.

The determinant of the related double-zeta matrix and its inverse have
published antecedents \cite{Zagier,Ma,MaTeo}. The central-factorial
construction is presented as a direct proof and an arithmetic extension
in these coordinates; exhaustive worldwide priority for its consequences
has not been established. A repository integration should preserve this
attribution instead of marking the known determinant as a newly settled
open conjecture.

\subsection{The mixed-color obstruction remains}

The weight-seven $S_6$ identity remains unresolved. The existing
coordinate change between $(g_{6,1},g_{4,3},g_{2,5})$ and
$(g_{6,1},g_{5,2},g_{4,3})$ uses two shuffle rows and is already proved
in the source. Its restricted-presentation obstructions likewise have
their own exact meanings. Integral diagonalization controls the
single-color doubles; it does not supply the missing mixed-color
functional equation.

A full weight-seven imaginary-word presentation would have 24514
admissible conjugation coordinates in the scoped convention. No
completed rank computation or reduction certificate for that enlarged
presentation is asserted here. This size is a planning datum, not an
independence result or an explanation that the conjecture must be true.

The retained $S_{10}$ and $S_{12}$ relations are also conjectures.
Their exact proximity intervals contain zero, so they cannot decide
equality. The rejected first $S_{12}$ vector, by contrast, has an interval
strictly separated from zero and is rigorously false. The archive keeps
the rejected search record for reproducibility and identifies the
retained vectors in a separate canonical input file.
